\documentclass[a4paper,12pt,openany]{article}


\usepackage{parskip}
\usepackage{graphicx} % Required for inserting images



\usepackage[pdftex]{graphicx}

\usepackage[utf8]{inputenc}
\usepackage[T1]{fontenc}
\usepackage[british]{babel}

\usepackage{amsmath}
\usepackage{amsthm}

\usepackage{typearea}

\usepackage[toc,page]{appendix}
\usepackage{float}
% Acronyms
\usepackage{arydshln}





% Divide your \newtheorem commands into groups
% Comment out the necessary commands

\theoremstyle{plain}
\newtheorem{theorem}{Theorem}
%\newtheorem{lemma}{Lemma}
\newtheorem{corollary}{Corollary}
\newtheorem{proposition}{Proposition}

\theoremstyle{definition}
\newtheorem{definition}{Definition}
%\newtheorem{condition}{Condition}
%\newtheorem{problem}{Problem}
%\newtheorem{example}{Example}

\theoremstyle{remark}
%\newtheorem{remark}{Remark}
%\newtheorem{note}{Note}
%\newtheorem{claim}{Claim}

\usepackage{newtxtext}
\usepackage{newtxmath}


\usepackage{amsmath}

\begin{document}


\begin{center}
\huge Algebraic Geometry\\
\vspace{0.2cm}
\Large Solutions to selected problems for assignment 3 \\  \vspace{0.5cm}
Mara Kali\v{c}anin Dimitrov
\vspace{0.5cm}
\end{center}



\subsection*{Problem 8.1}

Let $P=[a_0:\cdots:a_n]\in P^n$. Find a homogeneous ideal $I\subset k[x_0,\dots,x_n]$
such that
\[
V(I)=\{P\}.
\]

\subsection*{Restatement of the problem}

We are given a point
\[
P=[a_0:\cdots:a_n]\in P^n,
\]
and we must produce an explicit homogeneous ideal in the graded ring
$k[x_0,\dots,x_n]$ whose projective zero set is exactly the single point $P$.

\subsection*{Relevant definitions and results}

\begin{itemize}
    \item By Definition 8.1 (Projective $n$-space, p.~117), a point of $P^n$ is represented
    by a nonzero tuple $(a_0,\dots,a_n)$, where two such tuples define the same point if
    and only if they differ by multiplication by a nonzero scalar.

    \item By Definition 8.3 (p.~119), a homogeneous ideal is an ideal generated by
    homogeneous elements.

    \item By Definition 8.4 (zero of a homogeneous polynomial, p.~120), if
    $f\in k[x_0,\dots,x_n]$ is homogeneous, then a point
    $Q=[b_0:\cdots:b_n]\in P^n$ is a zero of $f$ if
    \[
    f(b_0,\dots,b_n)=0.
    \]

    \item By Definition 8.5 (projective algebraic set, p.~120), if $S$ is a set of
    homogeneous polynomials, then
    \[
    V(S)=\{Q\in P^n \mid f(Q)=0 \text{ for all } f\in S\}.
    \]
\end{itemize}

\subsection*{Complete mathematical solution}

We define
\[
I=\bigl(a_jx_i-a_ix_j \mid 0\le i,j\le n\bigr)\subset k[x_0,\dots,x_n].
\]

We first verify that $I$ is homogeneous. Each generator
\[
a_jx_i-a_ix_j
\]
is homogeneous of degree $1$. Therefore, by Definition 8.3 (p.~119), the ideal $I$
generated by these elements is a homogeneous ideal.

We now prove that
\[
V(I)=\{P\}.
\]

\medskip

\noindent\textbf{Step 1: Show that $P\in V(I)$.}

Let
\[
g_{ij}=a_jx_i-a_ix_j.
\]
Then
\[
g_{ij}(a_0,\dots,a_n)=a_ja_i-a_ia_j=0.
\]
Hence every generator of $I$ vanishes at $P$. By Definitions 8.4 and 8.5 (p.~120),
this implies
\[
P\in V(I).
\]

\medskip

\noindent\textbf{Step 2: Show that every point of $V(I)$ is equal to $P$.}

Let
\[
Q=[b_0:\cdots:b_n]\in V(I).
\]
Since $P$ is a point of $P^n$, Definition 8.1 (p.~117) implies that not all
$a_0,\dots,a_n$ are zero. Choose an index $s$ such that
\[
a_s\neq 0.
\]

Because $Q\in V(I)$ and the polynomial
\[
a_sx_i-a_ix_s
\]
belongs to $I$ for every $i$, we have
\[
a_sb_i-a_ib_s=0
\qquad\text{for all } i=0,\dots,n.
\]
We claim that $b_s\neq 0$. Indeed, if $b_s=0$, then the above relation becomes
\[
a_sb_i=0
\qquad\text{for all } i.
\]
Since $a_s\neq 0$, it follows that $b_i=0$ for all $i$, which is impossible because
$Q$ is a point of $P^n$ and therefore is represented by a nonzero tuple.

Thus $b_s\neq 0$. Define
\[
\lambda=\frac{b_s}{a_s}\in k^\times.
\]
Then for every $i$,
\[
a_sb_i=a_ib_s=a_i(\lambda a_s).
\]
Since $a_s\neq 0$, we may divide by $a_s$ and obtain
\[
b_i=\lambda a_i
\qquad\text{for all } i=0,\dots,n.
\]
Therefore,
\[
(b_0,\dots,b_n)=\lambda(a_0,\dots,a_n).
\]
By Definition 8.1 (p.~117), this means that
\[
Q=[b_0:\cdots:b_n]=[a_0:\cdots:a_n]=P.
\]

We have proved that every point of $V(I)$ is equal to $P$, so
\[
V(I)\subseteq \{P\}.
\]

Since Step 1 gave $\{P\}\subseteq V(I)$, we conclude that
\[
V(I)=\{P\}.
\]

\subsection*{Logical closure}

Hence one may take
\[
I=\bigl(a_jx_i-a_ix_j \mid 0\le i,j\le n\bigr),
\]
which is a homogeneous ideal in $k[x_0,\dots,x_n]$ whose projective zero set is
exactly the point $P$.

This proves the claim.




\subsection*{Problem 8.5}

Let $V_1,\dots,V_n$ be hyperplanes in $P^n$. Show that
\[
\bigcap_{i=1}^n V_i \neq \emptyset.
\]

\subsection*{Restatement of the problem}

We are given $n$ hyperplanes in projective $n$-space $P^n$, and we must
prove that they have a common point.

Since the problem uses the word \emph{hyperplane}, we interpret each
$V_i$ in the standard sense used in this chapter: namely, $V_i$ is the zero
set of a nonzero homogeneous polynomial of degree $1$. This is consistent
with the coordinate hyperplanes introduced in \S8.1 (p.~118) and with
Definitions 8.4 and 8.5 (p.~120), which describe projective algebraic sets as
zero sets of homogeneous polynomials.

\subsection*{Relevant definitions and results}

\begin{itemize}
    \item By Definition 8.1 (Projective $n$-space, p.~117), a point of $P^n$
    is an equivalence class
    \[
    [a_0:\cdots:a_n]
    \]
    with $(a_0,\dots,a_n)\neq (0,\dots,0)$, where two nonzero tuples define
    the same point if and only if they differ by multiplication by a nonzero
    scalar. We will use this to construct the desired common point.

    \item In \S8.1 (p.~118), the coordinate hyperplanes
    \[
    H_i=\{[a_0:\cdots:a_n]\in P^n\mid a_i=0\}
    \]
    are introduced. This indicates how hyperplanes in projective space are
    described by homogeneous linear equations.

    \item By Definition 8.4 (zero of a homogeneous polynomial, p.~120), if
    $\ell\in k[x_0,\dots,x_n]$ is homogeneous and
    $P=[a_0:\cdots:a_n]\in P^n$, then $P$ is a zero of $\ell$ precisely when
    $\ell(a_0,\dots,a_n)=0$.

    \item By Definition 8.5 (projective algebraic set, p.~120), for a set of
    homogeneous polynomials $S$,
    \[
    V(S)=\{P\in P^n\mid f(P)=0 \text{ for all } f\in S\}.
    \]
    In particular, if $V_i=V(\ell_i)$ for a linear form $\ell_i$, then proving
    that a point lies in all $V_i$ amounts to proving that all the $\ell_i$
    vanish at that point.
\end{itemize}

\subsection*{Complete mathematical solution}

For each $i=1,\dots,n$, let
\[
\ell_i(x_0,\dots,x_n)=a_{i0}x_0+\cdots+a_{in}x_n
\]
be a nonzero homogeneous polynomial of degree $1$ such that
\[
V_i=V(\ell_i).
\]
Thus we must show that there exists a point
\[
P=[b_0:\cdots:b_n]\in P^n
\]
such that
\[
\ell_i(b_0,\dots,b_n)=0
\qquad\text{for all } i=1,\dots,n.
\]

Equivalently, we must show that the homogeneous linear system
\[
\sum_{j=0}^n a_{ij}x_j=0,
\qquad i=1,\dots,n,
\]
has a nonzero solution $(x_0,\dots,x_n)\in k^{n+1}$.

We prove this by induction on $n$.

\medskip

\noindent\textbf{Base case: $n=1$.}

We have one nonzero linear equation in two variables:
\[
a_{10}x_0+a_{11}x_1=0,
\]
where $(a_{10},a_{11})\neq (0,0)$.
Consider the vector
\[
(-a_{11},a_{10}).
\]
This vector is nonzero, since $a_{10}$ and $a_{11}$ are not both zero, and
\[
a_{10}(-a_{11})+a_{11}a_{10}=0.
\]
Hence there exists a nonzero solution.

\medskip

\noindent\textbf{Induction step.}

Assume the statement proved for $n-1$, i.e. assume that any homogeneous
system of $n-1$ linear equations in $n$ variables has a nonzero solution.
We consider a homogeneous system of $n$ linear equations in $n+1$
variables:
\[
\sum_{j=0}^n a_{ij}x_j=0,
\qquad i=1,\dots,n.
\]

There are two cases.

\smallskip

\noindent\emph{Case 1: $a_{in}=0$ for all $i=1,\dots,n$.}

Then the variable $x_n$ does not appear in any equation. Therefore the
vector
\[
(0,\dots,0,1)
\]
is a nonzero solution of the whole system.

\smallskip

\noindent\emph{Case 2: $a_{in}\neq 0$ for at least one index $i$.}

After reordering the equations if necessary, we may assume that
\[
a_{1n}\neq 0.
\]
The first equation is then
\[
a_{10}x_0+\cdots+a_{1,n-1}x_{n-1}+a_{1n}x_n=0,
\]
so we may solve for $x_n$:
\[
x_n=-a_{1n}^{-1}\sum_{j=0}^{n-1}a_{1j}x_j.
\]

Substituting this expression into the remaining equations
\[
\sum_{j=0}^n a_{ij}x_j=0,
\qquad i=2,\dots,n,
\]
we obtain
\[
\sum_{j=0}^{n-1}\left(a_{ij}-a_{in}a_{1n}^{-1}a_{1j}\right)x_j=0,
\qquad i=2,\dots,n.
\]
Thus we obtain a homogeneous system of $n-1$ linear equations in the
$n$ variables $x_0,\dots,x_{n-1}$.

By the induction hypothesis, there exists a nonzero solution
\[
(b_0,\dots,b_{n-1})\in k^n
\]
to this reduced system. Define
\[
b_n=-a_{1n}^{-1}\sum_{j=0}^{n-1}a_{1j}b_j.
\]
Then $(b_0,\dots,b_n)$ satisfies the first equation by construction, and it
satisfies the remaining equations because $(b_0,\dots,b_{n-1})$ satisfies the
substituted system. Since $(b_0,\dots,b_{n-1})$ is nonzero, the vector
$(b_0,\dots,b_n)$ is also nonzero.

Hence the original homogeneous system of $n$ equations in $n+1$ variables
has a nonzero solution.

This completes the induction.

\medskip

We now return to the hyperplanes $V_1,\dots,V_n$. By the argument above,
there exists a nonzero tuple $(b_0,\dots,b_n)\in k^{n+1}$ such that
\[
\ell_i(b_0,\dots,b_n)=0
\qquad\text{for all } i=1,\dots,n.
\]
Since $(b_0,\dots,b_n)\neq (0,\dots,0)$, Definition 8.1 (p.~117) gives a
point
\[
P=[b_0:\cdots:b_n]\in P^n.
\]
For each $i$, Definition 8.4 (p.~120) implies that $P$ is a zero of $\ell_i$,
because
\[
\ell_i(b_0,\dots,b_n)=0.
\]
Therefore
\[
P\in V(\ell_i)=V_i
\qquad\text{for all } i=1,\dots,n.
\]
Hence
\[
P\in \bigcap_{i=1}^n V_i.
\]
It follows that
\[
\bigcap_{i=1}^n V_i\neq\emptyset.
\]

\subsection*{Logical closure}

We have proved that any $n$ hyperplanes in $P^n$ have a common point.
Therefore
\[
\bigcap_{i=1}^n V_i\neq\emptyset.
\]

This proves the claim.



\subsection*{Problem 9.1}

Let
\[
N=\frac{d(d+3)}{2}
\]
and let $P_1,\dots,P_N$ be points in $P^2$. Show that there is at least one
curve of degree $d$ passing through all of the points $P_i$.

\subsection*{Restatement of the problem}

We are given
\[
N=\frac{d(d+3)}{2}
\]
points $P_1,\dots,P_N \in P^2$, and we must prove that there exists a
projective plane curve $C$ of degree $d$ such that
\[
P_i \in C \qquad \text{for all } i=1,\dots,N.
\]

Equivalently, we must show that there exists a nonzero homogeneous
polynomial of degree $d$ whose associated projective plane curve contains
all the given points.

\subsection*{Relevant definitions and results}

\begin{itemize}
    \item By Definition 9.1 (Projective plane curve, \S9.1, p.~131), a
    projective plane curve is an equivalence class of nonzero homogeneous
    polynomials in $k[x_0,x_1,x_2]$ under multiplication by a nonzero scalar.
    This is the notion of ``curve of degree $d$'' used in the problem.

    \item At the beginning of \S9.2 (pp.~133--134), if $A_d$ denotes the
    vector space of homogeneous polynomials of degree $d$ in
    $k[x_0,x_1,x_2]$, then after choosing the monomial basis
    $M_0,\dots,M_N$, where
    \[
    N=\frac{(d+1)(d+2)}{2}-1=\frac{d(d+3)}{2},
    \]
    the set $C_d$ of plane curves of degree $d$ is identified with the
    projective space $P^N$. We will use this to translate the problem into
    a statement about hyperplanes in projective space.

    \item By Lemma 9.6 (\S9.2, p.~134), for any point $P\in P^2$, the set
    of curves of degree $d$ that contain $P$ forms a hyperplane in
    $P^{d(d+3)/2}$. We will apply this once for each of the points
    $P_1,\dots,P_N$.

    \item By Problem 8.5 (\S8.11, p.~130), if $V_1,\dots,V_n$ are
    hyperplanes in $P^n$, then
    \[
    \bigcap_{i=1}^n V_i \neq \emptyset.
    \]
    We will apply this with $n=N$.
\end{itemize}

\subsection*{Complete mathematical solution}

Let
\[
A=k[x_0,x_1,x_2],
\]
and let $A_d$ be its $d$th graded component, i.e. the vector space of
homogeneous polynomials of degree $d$.

Choose the monomials of degree $d$
\[
M_0,\dots,M_N
\]
as a basis of $A_d$, where
\[
N=\frac{(d+1)(d+2)}{2}-1=\frac{d(d+3)}{2}.
\]
Then every nonzero element $f\in A_d$ can be written uniquely in the form
\[
f=a_0M_0+\cdots+a_NM_N.
\]

By Definition 9.1 (\S9.1, p.~131), two nonzero homogeneous polynomials
define the same projective plane curve if and only if they differ by
multiplication by a nonzero scalar. Therefore the set $C_d$ of projective
plane curves of degree $d$ is naturally identified with the set of nonzero
coefficient vectors
\[
(a_0,\dots,a_N)
\]
modulo multiplication by a nonzero scalar. This is precisely the projective
space $P^N$.

For each $i=1,\dots,N$, define
\[
V_i=\{\, C\in C_d \mid P_i\in C \,\}.
\]
Thus $V_i$ is the set of all degree $d$ curves passing through the point
$P_i$.

Now Lemma 9.6 (\S9.2, p.~134) applies. Since each $P_i$ is a point of
$P^2$, the lemma tells us that $V_i$ is a hyperplane in
\[
P^{d(d+3)/2}=P^N.
\]
Hence
\[
V_1,\dots,V_N
\]
are $N$ hyperplanes in $P^N$.

We may therefore apply Problem 8.5 (\S8.11, p.~130) with $n=N$. Its
hypotheses are satisfied because we are considering exactly $N$ hyperplanes
inside $P^N$. We conclude that
\[
\bigcap_{i=1}^N V_i \neq \emptyset.
\]

Choose a point
\[
C\in \bigcap_{i=1}^N V_i.
\]
Since $C\in V_i$ for every $i$, the definition of $V_i$ gives
\[
P_i\in C \qquad \text{for all } i=1,\dots,N.
\]
Thus $C$ is a projective plane curve of degree $d$ passing through all the
given points.

\subsection*{Logical closure}

We have proved that any collection of
\[
N=\frac{d(d+3)}{2}
\]
points in $P^2$ lies on at least one projective plane curve of degree $d$.

This proves the claim.


\subsection*{Problem 9.2}

Let $P_1,\dots,P_r$ be points in $P^2$ in general position.

\begin{enumerate}
    \item[(a)] Show that there is a unique line passing through $P_1$ and $P_2$.
    \item[(b)] Show that there is a unique conic passing through $P_1,\dots,P_5$.
    \item[(c)] Show that there is a unique cubic passing through $P_1,\dots,P_9$.
    \item[(d)] How many points in general position are required to specify a quartic?
\end{enumerate}

\subsection*{Restatement of the problem}

For each degree $d=1,2,3,4$, we consider the linear system of plane curves
of degree $d$ passing through a prescribed collection of simple points in
$P^2$.

We must show that:

\begin{itemize}
    \item two points in general position determine a unique line,
    \item five points in general position determine a unique conic,
    \item nine points in general position determine a unique cubic,
    \item and we must determine the number of points in general position
    needed to determine a quartic uniquely.
\end{itemize}

Here the phrase ``in general position'' is understood in the sense of the
paragraph immediately preceding Theorem 9.16 (\S9.2, p.~138), applied to
the relevant degree $d$ and to simple base points (that is, all multiplicities
equal to $1$).

\subsection*{Relevant definitions and results}

\begin{itemize}
    \item By Definition 8.1 (Projective $n$-space, pp.~117--118), $P^0$
    consists of exactly one point. We will use this when a linear system has
    dimension $0$.

    \item By Examples 9.2, 9.3 and 9.4 (\S9.2, p.~134), the set of lines,
    conics, cubics and quartics in $P^2$ is identified with the projective
    spaces
    \[
    P^2,\qquad P^5,\qquad P^9,\qquad P^{14},
    \]
    respectively. Thus curves of degree $d$ are parametrized by a
    projective space of dimension
    \[
    \frac{d(d+3)}{2}.
    \]

    \item By Definition 9.5 (Linear system of curves of degree $d$, p.~134),
    a linear system is a linear subvariety of the projective space of curves
    of degree $d$.

    \item By Definition 9.10 (\S9.2, p.~135), if
    \[
    D=P_1+\cdots+P_r
    \]
    is an effective zero-cycle, then $L_d(-D)$ denotes the linear system of
    plane curves of degree $d$ passing through the points
    $P_1,\dots,P_r$.

    \item By Definition 9.12 (Virtual dimension and expected dimension,
    pp.~135--136), for
    \[
    D=P_1+\cdots+P_r
    \]
    we have
    \[
    \nu(L_d(-D))=\frac{d(d+3)}{2}-r
    \]
    and
    \[
    e(L_d(-D))=\max\{-1,\nu(L_d(-D))\}.
    \]
    We also use the convention stated there that
    \[
    \dim(\emptyset)=-1.
    \]

    \item By the paragraph immediately preceding Theorem 9.16 (p.~138),
    points are in general position precisely when the relevant evaluation map
    has maximal rank.

    \item By Theorem 9.16 (\S9.2, p.~138), if
    \[
    D=P_1+\cdots+P_r
    \]
    and the points are in general position, then
    \[
    \dim(L_d(-D))=e(L_d(-D)).
    \]
    This is the key tool for computing the dimensions of the relevant
    linear systems.
\end{itemize}

\subsection*{Complete mathematical solution}

We first make a uniform observation.

Let
\[
D_r=P_1+\cdots+P_r.
\]
By Definition 9.10 (\S9.2, p.~135), the linear system $L_d(-D_r)$ is
exactly the set of plane curves of degree $d$ passing through the points
$P_1,\dots,P_r$.

Since all multiplicities are equal to $1$, Definition 9.12
(pp.~135--136) gives
\[
\nu(L_d(-D_r))=\frac{d(d+3)}{2}-r.
\]
If the points are in general position, then Theorem 9.16 (p.~138) implies
\[
\dim(L_d(-D_r))=e(L_d(-D_r)).
\]

Therefore, whenever
\[
r=\frac{d(d+3)}{2},
\]
we get
\[
\nu(L_d(-D_r))=0,
\qquad
e(L_d(-D_r))=0,
\qquad
\dim(L_d(-D_r))=0.
\]

Now $L_d(-D_r)$ is a linear system, hence by Definition 9.5 (p.~134) a
linear subvariety of the projective space parametrizing degree $d$ curves.
A linear subvariety of dimension $0$ is a copy of $P^0$, and by
Definition 8.1 (pp.~117--118), $P^0$ consists of exactly one point.
Consequently, a linear system of dimension $0$ contains exactly one curve.

We now apply this observation in each case.

\medskip

\noindent\textbf{(a) Two points determine a unique line.}

For lines we have $d=1$. Then
\[
\frac{d(d+3)}{2}=\frac{1\cdot 4}{2}=2.
\]
Thus for
\[
D_2=P_1+P_2
\]
we obtain
\[
\dim(L_1(-D_2))=0
\]
by Theorem 9.16 and Definition 9.12.

Therefore the set of lines passing through $P_1$ and $P_2$ consists of
exactly one line.

This proves part \textup{(a)}.

\medskip

\noindent\textbf{(b) Five points determine a unique conic.}

For conics we have $d=2$. Then
\[
\frac{d(d+3)}{2}=\frac{2\cdot 5}{2}=5.
\]
Thus for
\[
D_5=P_1+\cdots+P_5
\]
we obtain
\[
\dim(L_2(-D_5))=0.
\]

Hence the linear system of conics passing through
$P_1,\dots,P_5$ consists of exactly one conic.

This proves part \textup{(b)}.

\medskip

\noindent\textbf{(c) Nine points determine a unique cubic.}

For cubics we have $d=3$. Then
\[
\frac{d(d+3)}{2}=\frac{3\cdot 6}{2}=9.
\]
Thus for
\[
D_9=P_1+\cdots+P_9
\]
we obtain
\[
\dim(L_3(-D_9))=0.
\]

Hence there is exactly one cubic passing through
$P_1,\dots,P_9$.

This proves part \textup{(c)}.

\medskip

\noindent\textbf{(d) Number of general points required to specify a quartic.}

For quartics we have $d=4$. By Example 9.4 (\S9.2, p.~134), quartics are
parametrized by $P^{14}$, and indeed
\[
\frac{d(d+3)}{2}=\frac{4\cdot 7}{2}=14.
\]

Let
\[
D_r=P_1+\cdots+P_r
\]
for $r$ points in general position. Then Theorem 9.16 and Definition 9.12
give
\[
\dim(L_4(-D_r))=e(L_4(-D_r)).
\]
If $r\le 14$, then
\[
e(L_4(-D_r))=14-r.
\]

For a quartic to be \emph{specified}, it must be determined uniquely.
As explained above, this happens exactly when the corresponding linear
system has dimension $0$. Thus we require
\[
14-r=0,
\]
so
\[
r=14.
\]

We also verify that this is the correct number:

\begin{itemize}
    \item if $r=14$, then $\dim(L_4(-D_{14}))=0$, so there is a unique
    quartic through the $14$ points;

    \item if $r<14$, then $\dim(L_4(-D_r))=14-r>0$, so there is a
    positive-dimensional family of quartics through the points, and hence
    the quartic is not uniquely determined.
\end{itemize}

Therefore, \emph{14 points in general position are required to specify a
quartic.}

This proves part \textup{(d)}.

\subsection*{Logical closure}

We have shown that:
\[
2 \text{ general points determine a unique line,}
\]
\[
5 \text{ general points determine a unique conic,}
\]
\[
9 \text{ general points determine a unique cubic,}
\]
and
\[
14 \text{ general points determine a unique quartic.}
\]

This completes the solution of Problem 9.2.

\end{document}

